\input{../tex-inputs/latex-korrekturansicht-vorspann}

\section[Arthur Schnitzler an Theodor Herzl, 25. 7. 1892]{L03900 Arthur Schnitzler an Theodor Herzl, 25. 7. 1892}
\nopagebreak\mylabel{L03900v}
\rehead{ }\normalsize\beginnumbering\briefempfaengerindex{Herzl, Theodor@\textsc{Herzl, Theodor}!zzzSchnitzler, Arthur@\emph{von Arthur Schnitzler}!1892-07-251@{25. 7. 1892}|(be}
\toendnotes[C]{\smallbreak\pagebreak[2]}
\correspDesc{Versand  durch Arthur Schnitzler am 25. 7. 1892 in Wien
\newline{}Erhalt  durch Theodor Herzl im Zeitraum [25. 7. 1892 – 28. 7. 1892?] in Paris}\toendnotes[C]{\smallbreak}
\Standort{Jerusalem, Central Zionist Archives, H1:1924-1.}
\physDesc{Brief, 1 Blatt, 3 Seiten, 719 Zeichen
\newline{}Handschrift: schwarze Tinte, deutsche Kurrent
\newline{}Ordnung: mit Bleistift von unbekannter Hand nummeriert: »(1)« und innerhalb das Konvoluts paginiert: »1«–»3« }
\buchAbdrucke{\weitereDrucke{Arthur Schnitzler: \emph{Briefe 1875–1912}. Herausgegeben von Therese Nickl und Heinrich Schnitzler. Frankfurt am Main: \emph{S. Fischer} 1981, S. 123.} }\toendnotes[C]{\smallbreak}
\pstart
           \raggedleft{}{\pb}25. 7. 92.\pend
           
\pstart
           \raggedleft{}\textcolor{pink}{Wien}\oindex{Wien@\textbf{Wien}, \emph{Verwaltungsgebiet}|pw}{}\ledrightnote{\textcolor{pink}{Wien}}.\pend
           
\pstart{}Verehrteſter Herr Doktor,\pend\vspace{0.5em}
\pstart
           es hat mich ganz beſonders gefreut, mit dem \textcolor{green}{Mährchen}\pwindex{Schnitzler, Arthur 15. 5. 1862 Wien – 21. 10. 1931 ebd.@\textsc{Schnitzler, Arthur} (15. 5. 1862 Wien – 21. 10. 1931 ebd.), \emph{Schriftsteller, Mediziner}!Märchen. Schauspiel in drei Aufzügen@\strich\emph{Das Märchen. Schauspiel in drei Aufzügen}|pw}{}\ledrightnote{\textcolor{green}{Das Märchen. Schauspiel in drei Aufzügen}} – wie mir \label{K_L03900-1v}\edtext{\textcolor{blue}{Goldma{\geminationn}}\pwindex{Goldmann, Paul 31.\,1.\,1865 Breslau – 25.\,9.\,1935 Wien@\textsc{Goldmann, Paul} (31.\,1.\,1865 Breslau – 25.\,9.\,1935 Wien), \emph{Schriftsteller, Journalist}|pw}{}\ledrightnote{\textcolor{blue}{Paul Goldmann}}
                  mittheilte}{\lemma{\textnormal{\emph{Goldmann
                  mittheilte}}}\Cendnote{\textnormal{Siehe Paul Goldmann an Arthur Schnitzler, 24. 6. [1892]. }}}\label{K_L03900-1} – vor Ihrem
               kunſtverſtändigen Urtheile Beifall gefunden zu haben. Was nun meine andre Sachen
               betrifft, die Sie zu leſen wünſchen, ſo thut es mir recht leid, Ihrem liebens{\pb}würdigen Intereſſe nur zum geringen Theil entſprechen zu
               können. Das meiſte iſt in Zeitungen minderer Verbreitung erſchienen; und nur \textcolor{green}{weniges}\pwindex{Schnitzler, Arthur 15. 5. 1862 Wien – 21. 10. 1931 ebd.@\textsc{Schnitzler, Arthur} (15. 5. 1862 Wien – 21. 10. 1931 ebd.), \emph{Schriftsteller, Mediziner}!Episode@\strich\emph{Episode}|pwv}\pwindex{Schnitzler, Arthur 15. 5. 1862 Wien – 21. 10. 1931 ebd.@\textsc{Schnitzler, Arthur} (15. 5. 1862 Wien – 21. 10. 1931 ebd.), \emph{Schriftsteller, Mediziner}!Reichtum. Erzählung@\strich\emph{Reichtum. Erzählung}|pwv}\pwindex{Schnitzler, Arthur 15. 5. 1862 Wien – 21. 10. 1931 ebd.@\textsc{Schnitzler, Arthur} (15. 5. 1862 Wien – 21. 10. 1931 ebd.), \emph{Schriftsteller, Mediziner}!Alkandi’s Lied@\strich\emph{Alkandi’s Lied}|pwv}{}\ledrightnote{{$\rightarrow$}\emph{\textcolor{green}{Episode}}{\newline}{$\rightarrow$}\emph{\textcolor{green}{Reichtum. Erzählung}}{\newline}{$\rightarrow$}\emph{\textcolor{green}{Alkandi’s Lied}}} ſteht mir in
               \label{K_L03900-2v}\edtext{Separatabdrucken}{\lemma{\textnormal{\emph{Separatabdrucken}}}\Cendnote{\textnormal{Dem Autor zur Verfügung gestellte, weitere
                  Drucke aus der Zeitung oder Zeitschrift, die nur den Text enthielten.}}}\label{K_L03900-2} zur Verfügung. Dieſes \textcolor{green}{wenige}\pwindex{Schnitzler, Arthur 15. 5. 1862 Wien – 21. 10. 1931 ebd.@\textsc{Schnitzler, Arthur} (15. 5. 1862 Wien – 21. 10. 1931 ebd.), \emph{Schriftsteller, Mediziner}!Episode@\strich\emph{Episode}|pwv}\pwindex{Schnitzler, Arthur 15. 5. 1862 Wien – 21. 10. 1931 ebd.@\textsc{Schnitzler, Arthur} (15. 5. 1862 Wien – 21. 10. 1931 ebd.), \emph{Schriftsteller, Mediziner}!Reichtum. Erzählung@\strich\emph{Reichtum. Erzählung}|pwv}{}\ledrightnote{{$\rightarrow$}\emph{\textcolor{green}{Episode}}{\newline}{$\rightarrow$}\emph{\textcolor{green}{Reichtum. Erzählung}}} ſchicke ich Ihnen hiemit; will Ihnen aber nicht
               verſchweigen, daſs im Herbſt {\pb}eine ganze Sa{\geminationm}lung von \textcolor{green}{\textsc{Anatol}ſtückeln}\pwindex{Schnitzler, Arthur 15. 5. 1862 Wien – 21. 10. 1931 ebd.@\textsc{Schnitzler, Arthur} (15. 5. 1862 Wien – 21. 10. 1931 ebd.), \emph{Schriftsteller, Mediziner}!Anatol@\strich\emph{Anatol}|pw}{}\ledrightnote{\textcolor{green}{Anatol}}\footnote{\noindent{}1. 
                     \textsc{\textcolor{green}{Episode}\pwindex{Schnitzler, Arthur 15. 5. 1862 Wien – 21. 10. 1931 ebd.@\textsc{Schnitzler, Arthur} (15. 5. 1862 Wien – 21. 10. 1931 ebd.), \emph{Schriftsteller, Mediziner}!Episode@\strich\emph{Episode}|pw}}} – u hoffentlich bald drauf noch was \textcolor{green}{andres}\pwindex{Schnitzler, Arthur 15. 5. 1862 Wien – 21. 10. 1931 ebd.@\textsc{Schnitzler, Arthur} (15. 5. 1862 Wien – 21. 10. 1931 ebd.), \emph{Schriftsteller, Mediziner}!Sterben. Novelle@\strich\emph{Sterben. Novelle}|pwuv}{}\ledrightnote{{$\rightarrow$}\emph{\textcolor{green}{Sterben. Novelle}}} von mir erſcheinen wird. –\pend
           
\pstart
           In aufrichtiger Verehrung{\\[\baselineskip]}Ihr ſehr ergebner{\\[\baselineskip]}\spacefill\mbox{D\textsuperscript{r}ArthSchnitzler}\pend
           \leftskip=0em{}\selectlanguage{ngerman}\endnumbering\briefempfaengerindex{Herzl, Theodor@\textsc{Herzl, Theodor}!zzzSchnitzler, Arthur@\emph{von Arthur Schnitzler}!1892-07-251@{25. 7. 1892}|)be}\mylabel{L03900h}  \newcommand{\dateiname}{L03900}\newcommand{\titel}{Arthur Schnitzler an Theodor Herzl, 25. 7. 1892}\newcommand{\editorInnen}{Herausgegeben von Jahnke, SelmaMüller, Martin Anton}\input{../tex-inputs/latex-korrekturansicht-abspann}
